
\documentclass{article}
\usepackage{colortbl}
\usepackage{makecell}
\usepackage{multirow}
\usepackage{supertabular}

\begin{document}

\newcounter{utterance}

\centering \large Interaction Transcript for game `function\_detective', experiment `list\_active\_inputs', episode 114 with Qwen3.6{-}35B{-}A3B{-}FP8{-}without{-}reasoning.
\vspace{24pt}

{ \footnotesize  \setcounter{utterance}{1}
\setlength{\tabcolsep}{0pt}
\begin{supertabular}{c@{$\;$}|p{.15\linewidth}@{}p{.15\linewidth}p{.15\linewidth}p{.15\linewidth}p{.15\linewidth}p{.15\linewidth}}
    \# & \multicolumn{2}{c}{Player} && \multicolumn{2}{c}{Game Master} \\
    \hline

    \theutterance \stepcounter{utterance}  
    & & & \multicolumn{4}{p{0.6\linewidth}}{
        \cellcolor[rgb]{0.9,0.9,0.9}{
            \makecell[{{p{\linewidth}}}]{
                \texttt{\tiny{[P1$\langle$GM]}}
                \texttt{I have a hidden Python function for you to figure out.} \\
\\ 
\texttt{Your job is to infer what it does and then submit a correct Python implementation.} \\
\\ 
\texttt{Function signature:} \\
\texttt{(items: List{[}int{]}) {-}> int} \\
\\ 
\texttt{You have 15 turns total.} \\
\\ 
\texttt{You choose test inputs x, and I return the exact output f(x).} \\
\\ 
\texttt{You must make your best SOLVE attempt no later than the second{-}to{-}last turn.} \\
\texttt{Do not wait until the very end.} \\
\\ 
\texttt{You can use these actions:} \\
\texttt{{-} JSON format described at the end of the prompt} \\
\texttt{{-} SOLVE: ```python ... ```} \\
\\ 
\texttt{SOLVE: ```python} \\
\texttt{def solution(<value1>):} \\
\texttt{    ...} \\
\texttt{```} \\
\\ 
\texttt{Do not add import statements in your guess.} \\
\texttt{Anything you might need is already available, so just write the function body and return the answer.} \\
\\ 
\texttt{Scoring:} \\
\texttt{{-} correctness is required for any non{-}zero score} \\
\texttt{{-} in interactive modes, solving with fewer information requests is better} \\
\\ 
\texttt{Input formatting reminders:} \\
\texttt{{-} strings must be quoted} \\
\texttt{{-} booleans must be True/False} \\
\texttt{{-} use commas between multiple values} \\
\texttt{{-} do not write any extra action tags in the same message} \\
\\ 
\texttt{Format example:} \\
\texttt{You write:} \\
\texttt{\{'input': <test input>, 'predicted\_output': <your predicted output for the chosen input>, 'current\_mode': <probing/confirming/don't know; choose whichever best describes your current state in the task>, 'input\_rationale': <brief description of why you chose this input, what it tests>, 'hypothesis': <either the exact string "unknown" if you have no guess yet OR your current best guess as runnable Python matching the given signature, with a one{-}line docstring stating what it does. Never write placeholder code, if you cannot commit to a guess, use "unknown". Match the given signature.} \\
\texttt{e.g.} \\
\texttt{def f(x):} \\
\texttt{	"""Returns the absolute value of x."""} \\
\texttt{	return x if x >= 0 else {-}x} \\
\\ 
\texttt{Use "unknown" if you have no guess yet.>, 'confidence\_score': <score between 1{-}100. 100 means highest confidence in your hypothesis>\}} \\
\\ 
\texttt{I answer:} \\
\texttt{OUTPUT: <output\_number>} \\
\\ 
\texttt{Output only one valid json per turn.} \\
            }
        }
    }
     \\ \\

    \theutterance \stepcounter{utterance}  
    & \multicolumn{4}{p{0.6\linewidth}}{
        \cellcolor[rgb]{0.561,0.737,0.561}{
            \makecell[{{p{\linewidth}}}]{
                \texttt{\tiny{[P1$\rangle$GM]}}
                \texttt{To figure out what the hidden function does, I'll start by testing some simple inputs to understand its behavior. I'll begin with a few basic cases like an empty list, a single element, and a small list with distinct elements.} \\
\\ 
\texttt{\{'input': {[}{]}, 'predicted\_output': 0, 'current\_mode': 'probing', 'input\_rationale': 'Testing with an empty list to see the base case behavior.', 'hypothesis': 'unknown', 'confidence\_score': 10\}} \\
            }
        }
    }
    & & \\ \\

    \theutterance \stepcounter{utterance}  
    & & & \multicolumn{4}{p{0.6\linewidth}}{
        \cellcolor[rgb]{0.9,0.9,0.9}{
            \makecell[{{p{\linewidth}}}]{
                \texttt{\tiny{[P1$\langle$GM]}}
                \texttt{Game Violation: ParseError: Response is not valid JSON.} \\
\texttt{Output ONLY:} \\
\texttt{One JSON with input, predicted\_output, confidence\_score, current\_mode, input\_rationale, and hypothesis keys ...} \\
\texttt{OR} \\
\texttt{SOLVE: ```python ... ```} \\
            }
        }
    }
     \\ \\

    \theutterance \stepcounter{utterance}  
    & \multicolumn{4}{p{0.6\linewidth}}{
        \cellcolor[rgb]{0.561,0.737,0.561}{
            \makecell[{{p{\linewidth}}}]{
                \texttt{\tiny{[P1$\rangle$GM]}}
                \texttt{\{"input": {[}{]}, "predicted\_output": 0, "current\_mode": "probing", "input\_rationale": "Testing with an empty list to see the base case behavior.", "hypothesis": "unknown", "confidence\_score": 10\}} \\
            }
        }
    }
    & & \\ \\

    \theutterance \stepcounter{utterance}  
    & & & \multicolumn{2}{p{0.3\linewidth}}{
        \cellcolor[rgb]{0.9,0.9,0.9}{
            \makecell[{{p{\linewidth}}}]{
                \texttt{\tiny{[GM$|$GM]}}
                \texttt{Round: 1} \\
\texttt{Inputs: {[}{]}} \\
\texttt{Output: 0\#0} \\
            }
        }
    }
    & & \\ \\

    \theutterance \stepcounter{utterance}  
    & & & \multicolumn{4}{p{0.6\linewidth}}{
        \cellcolor[rgb]{0.9,0.9,0.9}{
            \makecell[{{p{\linewidth}}}]{
                \texttt{\tiny{[P1$\langle$GM]}}
                \texttt{OUTPUT: 0} \\
            }
        }
    }
     \\ \\

    \theutterance \stepcounter{utterance}  
    & \multicolumn{4}{p{0.6\linewidth}}{
        \cellcolor[rgb]{0.561,0.737,0.561}{
            \makecell[{{p{\linewidth}}}]{
                \texttt{\tiny{[P1$\rangle$GM]}}
                \texttt{The function returns 0 for an empty list. Let's test with a single element list to see if it returns the element itself or something else (like 0 or the element squared).} \\
\\ 
\texttt{\{"input": {[}5{]}, "predicted\_output": 5, "current\_mode": "probing", "input\_rationale": "Testing with a single element list to determine if the function returns the element itself.", "hypothesis": "unknown", "confidence\_score": 10\}} \\
            }
        }
    }
    & & \\ \\

    \theutterance \stepcounter{utterance}  
    & & & \multicolumn{4}{p{0.6\linewidth}}{
        \cellcolor[rgb]{0.9,0.9,0.9}{
            \makecell[{{p{\linewidth}}}]{
                \texttt{\tiny{[P1$\langle$GM]}}
                \texttt{Game Violation: ParseError: Response is not valid JSON.} \\
\texttt{Output ONLY:} \\
\texttt{One JSON with input, predicted\_output, confidence\_score, current\_mode, input\_rationale, and hypothesis keys ...} \\
\texttt{OR} \\
\texttt{SOLVE: ```python ... ```} \\
            }
        }
    }
     \\ \\

    \theutterance \stepcounter{utterance}  
    & \multicolumn{4}{p{0.6\linewidth}}{
        \cellcolor[rgb]{0.561,0.737,0.561}{
            \makecell[{{p{\linewidth}}}]{
                \texttt{\tiny{[P1$\rangle$GM]}}
                \texttt{\{"input": {[}5{]}, "predicted\_output": 5, "current\_mode": "probing", "input\_rationale": "Testing with a single element list to determine if the function returns the element itself.", "hypothesis": "unknown", "confidence\_score": 10\}} \\
            }
        }
    }
    & & \\ \\

    \theutterance \stepcounter{utterance}  
    & & & \multicolumn{2}{p{0.3\linewidth}}{
        \cellcolor[rgb]{0.9,0.9,0.9}{
            \makecell[{{p{\linewidth}}}]{
                \texttt{\tiny{[GM$|$GM]}}
                \texttt{Round: 2} \\
\texttt{Inputs: {[}5{]}} \\
\texttt{Output: 1\#1} \\
            }
        }
    }
    & & \\ \\

    \theutterance \stepcounter{utterance}  
    & & & \multicolumn{4}{p{0.6\linewidth}}{
        \cellcolor[rgb]{0.9,0.9,0.9}{
            \makecell[{{p{\linewidth}}}]{
                \texttt{\tiny{[P1$\langle$GM]}}
                \texttt{OUTPUT: 1} \\
            }
        }
    }
     \\ \\

    \theutterance \stepcounter{utterance}  
    & \multicolumn{4}{p{0.6\linewidth}}{
        \cellcolor[rgb]{0.561,0.737,0.561}{
            \makecell[{{p{\linewidth}}}]{
                \texttt{\tiny{[P1$\rangle$GM]}}
                \texttt{The function returns 1 for `{[}5{]}`. This is interesting. It's not the element itself. Let's try a few more cases to narrow down the possibilities. Maybe it's the length of the list? Or maybe it's related to the number of unique elements? Or perhaps it's a specific mathematical operation.} \\
\\ 
\texttt{Let's test `{[}1, 2{]}` and `{[}1, 1{]}` to distinguish between length{-}based and value{-}based logic.} \\
\\ 
\texttt{\{"input": {[}1, 2{]}, "predicted\_output": 2, "current\_mode": "probing", "input\_rationale": "Testing with two distinct elements to see if the output is the length (2) or something else.", "hypothesis": "unknown", "confidence\_score": 10\}} \\
            }
        }
    }
    & & \\ \\

    \theutterance \stepcounter{utterance}  
    & & & \multicolumn{4}{p{0.6\linewidth}}{
        \cellcolor[rgb]{0.9,0.9,0.9}{
            \makecell[{{p{\linewidth}}}]{
                \texttt{\tiny{[P1$\langle$GM]}}
                \texttt{Game Violation: ParseError: Response is not valid JSON.} \\
\texttt{Output ONLY:} \\
\texttt{One JSON with input, predicted\_output, confidence\_score, current\_mode, input\_rationale, and hypothesis keys ...} \\
\texttt{OR} \\
\texttt{SOLVE: ```python ... ```} \\
            }
        }
    }
     \\ \\

    \theutterance \stepcounter{utterance}  
    & \multicolumn{4}{p{0.6\linewidth}}{
        \cellcolor[rgb]{0.561,0.737,0.561}{
            \makecell[{{p{\linewidth}}}]{
                \texttt{\tiny{[P1$\rangle$GM]}}
                \texttt{\{"input": {[}1, 2{]}, "predicted\_output": 2, "current\_mode": "probing", "input\_rationale": "Testing with two distinct elements to see if the output is the length (2) or something else.", "hypothesis": "unknown", "confidence\_score": 10\}} \\
            }
        }
    }
    & & \\ \\

    \theutterance \stepcounter{utterance}  
    & & & \multicolumn{2}{p{0.3\linewidth}}{
        \cellcolor[rgb]{0.9,0.9,0.9}{
            \makecell[{{p{\linewidth}}}]{
                \texttt{\tiny{[GM$|$GM]}}
                \texttt{Round: 3} \\
\texttt{Inputs: {[}1, 2{]}} \\
\texttt{Output: 1\#1} \\
            }
        }
    }
    & & \\ \\

    \theutterance \stepcounter{utterance}  
    & & & \multicolumn{4}{p{0.6\linewidth}}{
        \cellcolor[rgb]{0.9,0.9,0.9}{
            \makecell[{{p{\linewidth}}}]{
                \texttt{\tiny{[P1$\langle$GM]}}
                \texttt{OUTPUT: 1} \\
            }
        }
    }
     \\ \\

    \theutterance \stepcounter{utterance}  
    & \multicolumn{4}{p{0.6\linewidth}}{
        \cellcolor[rgb]{0.561,0.737,0.561}{
            \makecell[{{p{\linewidth}}}]{
                \texttt{\tiny{[P1$\rangle$GM]}}
                \texttt{The function returns 1 for `{[}1, 2{]}`.} \\
\texttt{So far:} \\
\texttt{{-} `{[}{]}` {-}> 0} \\
\texttt{{-} `{[}5{]}` {-}> 1} \\
\texttt{{-} `{[}1, 2{]}` {-}> 1} \\
\\ 
\texttt{This suggests the output might be related to the number of unique elements, or perhaps something about the values themselves. Let's test `{[}1, 1{]}` to see if duplicates matter. If it's unique count, `{[}1, 1{]}` should be 1. If it's something else, it might differ.} \\
\\ 
\texttt{\{"input": {[}1, 1{]}, "predicted\_output": 1, "current\_mode": "probing", "input\_rationale": "Testing with duplicate elements to see if the function counts unique elements or total elements.", "hypothesis": "unknown", "confidence\_score": 10\}} \\
            }
        }
    }
    & & \\ \\

    \theutterance \stepcounter{utterance}  
    & & & \multicolumn{4}{p{0.6\linewidth}}{
        \cellcolor[rgb]{0.9,0.9,0.9}{
            \makecell[{{p{\linewidth}}}]{
                \texttt{\tiny{[P1$\langle$GM]}}
                \texttt{Game Violation: ParseError: Response is not valid JSON.} \\
\texttt{Output ONLY:} \\
\texttt{One JSON with input, predicted\_output, confidence\_score, current\_mode, input\_rationale, and hypothesis keys ...} \\
\texttt{OR} \\
\texttt{SOLVE: ```python ... ```} \\
            }
        }
    }
     \\ \\

    \theutterance \stepcounter{utterance}  
    & \multicolumn{4}{p{0.6\linewidth}}{
        \cellcolor[rgb]{0.561,0.737,0.561}{
            \makecell[{{p{\linewidth}}}]{
                \texttt{\tiny{[P1$\rangle$GM]}}
                \texttt{\{"input": {[}1, 1{]}, "predicted\_output": 1, "current\_mode": "probing", "input\_rationale": "Testing with duplicate elements to see if the function counts unique elements or total elements.", "hypothesis": "unknown", "confidence\_score": 10\}} \\
            }
        }
    }
    & & \\ \\

    \theutterance \stepcounter{utterance}  
    & & & \multicolumn{2}{p{0.3\linewidth}}{
        \cellcolor[rgb]{0.9,0.9,0.9}{
            \makecell[{{p{\linewidth}}}]{
                \texttt{\tiny{[GM$|$GM]}}
                \texttt{Round: 4} \\
\texttt{Inputs: {[}1, 1{]}} \\
\texttt{Output: 2\#2} \\
            }
        }
    }
    & & \\ \\

    \theutterance \stepcounter{utterance}  
    & & & \multicolumn{4}{p{0.6\linewidth}}{
        \cellcolor[rgb]{0.9,0.9,0.9}{
            \makecell[{{p{\linewidth}}}]{
                \texttt{\tiny{[P1$\langle$GM]}}
                \texttt{OUTPUT: 2} \\
            }
        }
    }
     \\ \\

    \theutterance \stepcounter{utterance}  
    & \multicolumn{4}{p{0.6\linewidth}}{
        \cellcolor[rgb]{0.561,0.737,0.561}{
            \makecell[{{p{\linewidth}}}]{
                \texttt{\tiny{[P1$\rangle$GM]}}
                \texttt{The function returns 2 for `{[}1, 1{]}`.} \\
\texttt{So far:} \\
\texttt{{-} `{[}{]}` {-}> 0} \\
\texttt{{-} `{[}5{]}` {-}> 1} \\
\texttt{{-} `{[}1, 2{]}` {-}> 1} \\
\texttt{{-} `{[}1, 1{]}` {-}> 2} \\
\\ 
\texttt{This is very interesting.} \\
\texttt{{-} `{[}5{]}` (length 1) {-}> 1} \\
\texttt{{-} `{[}1, 2{]}` (length 2) {-}> 1} \\
\texttt{{-} `{[}1, 1{]}` (length 2) {-}> 2} \\
\\ 
\texttt{It doesn't seem to be simply the length.} \\
\texttt{Let's look at the values.} \\
\texttt{{-} `{[}5{]}` {-}> 1. Maybe `5 \% 4`? No, that's 1.} \\
\texttt{{-} `{[}1, 2{]}` {-}> 1. Sum is 3. `3 \% 2`? No.} \\
\texttt{{-} `{[}1, 1{]}` {-}> 2. Sum is 2.} \\
\\ 
\texttt{Let's try a list with 3 elements to see if it scales.} \\
\texttt{`{[}1, 2, 3{]}` {-}> ?} \\
\texttt{`{[}1, 1, 1{]}` {-}> ?} \\
\\ 
\texttt{Let's test `{[}1, 2, 3{]}`. If it's unique count, it would be 3. If it's something else, we'll see.} \\
\\ 
\texttt{\{"input": {[}1, 2, 3{]}, "predicted\_output": 3} \\
            }
        }
    }
    & & \\ \\

    \theutterance \stepcounter{utterance}  
    & & & \multicolumn{4}{p{0.6\linewidth}}{
        \cellcolor[rgb]{0.9,0.9,0.9}{
            \makecell[{{p{\linewidth}}}]{
                \texttt{\tiny{[P1$\langle$GM]}}
                \texttt{Game Violation: ParseError: Response is not valid JSON.} \\
\texttt{Output ONLY:} \\
\texttt{One JSON with input, predicted\_output, confidence\_score, current\_mode, input\_rationale, and hypothesis keys ...} \\
\texttt{OR} \\
\texttt{SOLVE: ```python ... ```} \\
            }
        }
    }
     \\ \\

    \theutterance \stepcounter{utterance}  
    & \multicolumn{4}{p{0.6\linewidth}}{
        \cellcolor[rgb]{0.561,0.737,0.561}{
            \makecell[{{p{\linewidth}}}]{
                \texttt{\tiny{[P1$\rangle$GM]}}
                \texttt{\{"input": {[}1, 2, 3{]}, "predicted\_output": 3, "current\_mode": "probing", "input\_rationale": "Testing with three distinct elements to check if the output is the count of unique elements or something else.", "hypothesis": "unknown", "confidence\_score": 10\}} \\
            }
        }
    }
    & & \\ \\

    \theutterance \stepcounter{utterance}  
    & & & \multicolumn{2}{p{0.3\linewidth}}{
        \cellcolor[rgb]{0.9,0.9,0.9}{
            \makecell[{{p{\linewidth}}}]{
                \texttt{\tiny{[GM$|$GM]}}
                \texttt{Round: 5} \\
\texttt{Inputs: {[}1, 2, 3{]}} \\
\texttt{Output: 2\#2} \\
            }
        }
    }
    & & \\ \\

    \theutterance \stepcounter{utterance}  
    & & & \multicolumn{4}{p{0.6\linewidth}}{
        \cellcolor[rgb]{0.9,0.9,0.9}{
            \makecell[{{p{\linewidth}}}]{
                \texttt{\tiny{[P1$\langle$GM]}}
                \texttt{OUTPUT: 2} \\
            }
        }
    }
     \\ \\

    \theutterance \stepcounter{utterance}  
    & \multicolumn{4}{p{0.6\linewidth}}{
        \cellcolor[rgb]{0.561,0.737,0.561}{
            \makecell[{{p{\linewidth}}}]{
                \texttt{\tiny{[P1$\rangle$GM]}}
                \texttt{The function returns 2 for `{[}1, 2, 3{]}`.} \\
\texttt{So far:} \\
\texttt{{-} `{[}{]}` {-}> 0} \\
\texttt{{-} `{[}5{]}` {-}> 1} \\
\texttt{{-} `{[}1, 2{]}` {-}> 1} \\
\texttt{{-} `{[}1, 1{]}` {-}> 2} \\
\texttt{{-} `{[}1, 2, 3{]}` {-}> 2} \\
\\ 
\texttt{Let's analyze the pattern:} \\
\texttt{{-} Length 0 {-}> 0} \\
\texttt{{-} Length 1 {-}> 1} \\
\texttt{{-} Length 2 (distinct) {-}> 1} \\
\texttt{{-} Length 2 (same) {-}> 2} \\
\texttt{{-} Length 3 (distinct) {-}> 2} \\
\\ 
\texttt{This looks like it might be related to the number of unique elements, but `{[}1, 1{]}` having 1 unique element returning 2 contradicts that.} \\
\texttt{Wait, let's look at the sum of elements?} \\
\texttt{{-} `{[}{]}` sum 0 {-}> 0} \\
\texttt{{-} `{[}5{]}` sum 5 {-}> 1} \\
\texttt{{-} `{[}1, 2{]}` sum 3 {-}> 1} \\
\texttt{{-} `{[}1, 1{]}` sum 2 {-}> 2} \\
\texttt{{-} `{[}1, 2, 3{]}` sum 6 {-}> 2} \\
\\ 
\texttt{No obvious pattern with sum.} \\
\\ 
\texttt{Let's look at the maximum element?} \\
\texttt{{-} `{[}{]}` {-}> 0 (maybe default)} \\
\texttt{{-} `{[}5{]}` {-}> 5 != 1} \\
\texttt{{-} `{[}1, 2{]}` {-}>} \\
            }
        }
    }
    & & \\ \\

    \theutterance \stepcounter{utterance}  
    & & & \multicolumn{4}{p{0.6\linewidth}}{
        \cellcolor[rgb]{0.9,0.9,0.9}{
            \makecell[{{p{\linewidth}}}]{
                \texttt{\tiny{[P1$\langle$GM]}}
                \texttt{Game Violation: ParseError: Response is not valid JSON.} \\
\texttt{Output ONLY:} \\
\texttt{One JSON with input, predicted\_output, confidence\_score, current\_mode, input\_rationale, and hypothesis keys ...} \\
\texttt{OR} \\
\texttt{SOLVE: ```python ... ```} \\
            }
        }
    }
     \\ \\

    \theutterance \stepcounter{utterance}  
    & \multicolumn{4}{p{0.6\linewidth}}{
        \cellcolor[rgb]{0.561,0.737,0.561}{
            \makecell[{{p{\linewidth}}}]{
                \texttt{\tiny{[P1$\rangle$GM]}}
                \texttt{\{"input": {[}1, 2, 3, 4{]}, "predicted\_output": 2, "current\_mode": "probing", "input\_rationale": "Testing with four distinct elements to see if the output continues to be 2 or changes, helping to identify if it's related to unique counts or some other property.", "hypothesis": "unknown", "confidence\_score": 10\}} \\
            }
        }
    }
    & & \\ \\

    \theutterance \stepcounter{utterance}  
    & & & \multicolumn{2}{p{0.3\linewidth}}{
        \cellcolor[rgb]{0.9,0.9,0.9}{
            \makecell[{{p{\linewidth}}}]{
                \texttt{\tiny{[GM$|$GM]}}
                \texttt{Round: 6} \\
\texttt{Inputs: {[}1, 2, 3, 4{]}} \\
\texttt{Output: 2\#2} \\
            }
        }
    }
    & & \\ \\

    \theutterance \stepcounter{utterance}  
    & & & \multicolumn{4}{p{0.6\linewidth}}{
        \cellcolor[rgb]{0.9,0.9,0.9}{
            \makecell[{{p{\linewidth}}}]{
                \texttt{\tiny{[P1$\langle$GM]}}
                \texttt{OUTPUT: 2} \\
            }
        }
    }
     \\ \\

    \theutterance \stepcounter{utterance}  
    & \multicolumn{4}{p{0.6\linewidth}}{
        \cellcolor[rgb]{0.561,0.737,0.561}{
            \makecell[{{p{\linewidth}}}]{
                \texttt{\tiny{[P1$\rangle$GM]}}
                \texttt{\{"input": {[}1, 1, 1, 1{]}, "predicted\_output": 4, "current\_mode": "probing", "input\_rationale": "Testing with four identical elements to see if the output is the length of the list when all elements are the same.", "hypothesis": "unknown", "confidence\_score": 10\}} \\
            }
        }
    }
    & & \\ \\

    \theutterance \stepcounter{utterance}  
    & & & \multicolumn{2}{p{0.3\linewidth}}{
        \cellcolor[rgb]{0.9,0.9,0.9}{
            \makecell[{{p{\linewidth}}}]{
                \texttt{\tiny{[GM$|$GM]}}
                \texttt{Round: 7} \\
\texttt{Inputs: {[}1, 1, 1, 1{]}} \\
\texttt{Output: 4\#4} \\
            }
        }
    }
    & & \\ \\

    \theutterance \stepcounter{utterance}  
    & & & \multicolumn{4}{p{0.6\linewidth}}{
        \cellcolor[rgb]{0.9,0.9,0.9}{
            \makecell[{{p{\linewidth}}}]{
                \texttt{\tiny{[P1$\langle$GM]}}
                \texttt{OUTPUT: 4} \\
            }
        }
    }
     \\ \\

    \theutterance \stepcounter{utterance}  
    & \multicolumn{4}{p{0.6\linewidth}}{
        \cellcolor[rgb]{0.561,0.737,0.561}{
            \makecell[{{p{\linewidth}}}]{
                \texttt{\tiny{[P1$\rangle$GM]}}
                \texttt{\{"input": {[}1, 2, 2, 3{]}, "predicted\_output": 3, "current\_mode": "probing", "input\_rationale": "Testing with a mix of distinct and duplicate elements. If the function counts unique elements, it should be 3. If it counts total elements, it should be 4. If it's something else, we'll see.", "hypothesis": "unknown", "confidence\_score": 10\}} \\
            }
        }
    }
    & & \\ \\

    \theutterance \stepcounter{utterance}  
    & & & \multicolumn{2}{p{0.3\linewidth}}{
        \cellcolor[rgb]{0.9,0.9,0.9}{
            \makecell[{{p{\linewidth}}}]{
                \texttt{\tiny{[GM$|$GM]}}
                \texttt{Round: 8} \\
\texttt{Inputs: {[}1, 2, 2, 3{]}} \\
\texttt{Output: 2\#2} \\
            }
        }
    }
    & & \\ \\

    \theutterance \stepcounter{utterance}  
    & & & \multicolumn{4}{p{0.6\linewidth}}{
        \cellcolor[rgb]{0.9,0.9,0.9}{
            \makecell[{{p{\linewidth}}}]{
                \texttt{\tiny{[P1$\langle$GM]}}
                \texttt{OUTPUT: 2} \\
            }
        }
    }
     \\ \\

    \theutterance \stepcounter{utterance}  
    & \multicolumn{4}{p{0.6\linewidth}}{
        \cellcolor[rgb]{0.561,0.737,0.561}{
            \makecell[{{p{\linewidth}}}]{
                \texttt{\tiny{[P1$\rangle$GM]}}
                \texttt{\{"input": {[}1, 2, 2, 3{]}, "predicted\_output": 2, "current\_mode": "probing", "input\_rationale": "Testing with a mix of distinct and duplicate elements. The output is 2, which is the number of unique elements (1, 2, 3 {-}> 3 unique? No, 1, 2, 3 are 3 unique. Wait. 1, 2, 2, 3 has unique elements \{1, 2, 3\}, count is 3. But output is 2. So it's not unique count.", "hypothesis": "unknown", "confidence\_score": 10\}} \\
            }
        }
    }
    & & \\ \\

    \theutterance \stepcounter{utterance}  
    & & & \multicolumn{2}{p{0.3\linewidth}}{
        \cellcolor[rgb]{0.9,0.9,0.9}{
            \makecell[{{p{\linewidth}}}]{
                \texttt{\tiny{[GM$|$GM]}}
                \texttt{Round: 9} \\
\texttt{Inputs: {[}1, 2, 2, 3{]}} \\
\texttt{Output: 2\#2} \\
            }
        }
    }
    & & \\ \\

    \theutterance \stepcounter{utterance}  
    & & & \multicolumn{2}{p{0.3\linewidth}}{
        \cellcolor[rgb]{0.9,0.9,0.9}{
            \makecell[{{p{\linewidth}}}]{
                \texttt{\tiny{[GM$|$GM]}}
                \texttt{Game Verdict: LOSS (turn limit reached)} \\
            }
        }
    }
    & & \\ \\

\end{supertabular}
}

\end{document}
